\documentclass[11pt]{amsart}
\usepackage[margin=1.1in]{geometry}
\usepackage{amsmath,amssymb,amsthm}
\usepackage{booktabs}
\usepackage[colorlinks=true,linkcolor=blue,citecolor=blue,urlcolor=blue]{hyperref}

\newtheorem{theorem}{Theorem}
\newtheorem{proposition}[theorem]{Proposition}
\newtheorem{lemma}[theorem]{Lemma}
\newtheorem{corollary}[theorem]{Corollary}
\theoremstyle{definition}
\newtheorem{definition}[theorem]{Definition}
\newtheorem{problem}[theorem]{Problem}
\theoremstyle{remark}
\newtheorem{remark}[theorem]{Remark}

\newcommand{\lbl}[1]{\textup{\textsc{#1}}}
\newcommand{\DHL}{\mathrm{DHL}}
\newcommand{\R}{\mathcal{R}}

\title[Toward the twin prime conjecture]{Toward the twin prime conjecture: certificates, limits of the sieve, and a working programme}
\author{Jerry Kim}
\thanks{Prepared with the assistance of Claude (Anthropic). The computations, certificates and Lean sources accompany this paper.}
\date{Version 0.3, 12 November 2026}

\begin{document}

\begin{abstract}
This is a living working paper. It does \emph{not} prove the twin prime
conjecture, and it does not improve the best known bound on small gaps between
primes. It records (i) the state of the problem as of the date above, (ii) what
this project has verified for itself, by exact computation and by the Lean~4
kernel, (iii) small results of our own: a no-go statement for a natural
attempt to sharpen the pigeonhole step, on the plain and on the
$\varepsilon$-enlarged simplex; a generating-function formula that makes
certified computation with all symmetric test functions cheap; and an upper
bound showing that the $\varepsilon$-enlargement of Polymath8b cannot, by
itself, reach any $k\le29$, and (iv) a programme of original
directions that later versions will test. Every statement carries one of the
labels \lbl{Lean-checked}, \lbl{Computed}, \lbl{Proved here},
\lbl{Literature}, \lbl{Claimed} (a preprint not yet independently
verified), or \lbl{Open}.
\end{abstract}

\maketitle

\section{The problem and the honest position}

Write $p_n$ for the $n$-th prime and $H_1=\liminf_{n\to\infty}(p_{n+1}-p_n)$.
The twin prime conjecture is the statement $H_1=2$. It is \lbl{Open}.

A finite set $\mathcal H=\{h_1<\dots<h_k\}$ of integers is \emph{admissible} if
for every prime $p$ it misses at least one residue class modulo $p$. Let
$\DHL[k,2]$ be the statement that every admissible $k$-tuple has infinitely
many translates $n+\mathcal H$ containing at least two primes, and let $H(k)$
be the least diameter $h_k-h_1$ of an admissible $k$-tuple. Then
\begin{equation}\label{eq:reduction}
\DHL[k,2]\ \Longrightarrow\ H_1\le H(k),
\end{equation}
and $\DHL[2,2]$ applied to $\{0,2\}$ is the twin prime conjecture. All progress
on $H_1$ since 2013 consists of proving $\DHL[k,2]$ for smaller $k$.

\subsection*{Where the bound stands (12 November 2026)}

\begin{center}
\begin{tabular}{@{}rlll@{}}
\toprule
$H_1\le$ & $k$ & Source & Label\\
\midrule
$7\times10^7$ & $3.5\times10^6$ & Zhang \cite{Zha14} & \lbl{Literature}\\
600 & 105 & Maynard \cite{May15} & \lbl{Literature}\\
246 & 50 & Polymath8b \cite{Pol14b} & \lbl{Literature}\\
240 & 49$^{\ast}$ & Stadlmann \cite{Sta26} & \lbl{Claimed}\\
236 & 48$^{\ast}$ & Kintali \cite{Kin26}; Song--Yue \cite{SY26} & \lbl{Claimed}\\
212 & 45 & Charton et al.\ \cite{Cha26} & \lbl{Claimed}\\
186 & 40 & OpenAI \cite{OAI26} & \lbl{Claimed}\\
182 & 39 & GPT-6 Astra, prompted by Sharifi \cite{Sha26} & \lbl{Claimed}\\
\midrule
12 & 5 & Maynard, assuming EH \cite{May15} & \lbl{Literature}\\
6 & 3 & Polymath8b, assuming GEH \cite{Pol14b} & \lbl{Literature}\\
\bottomrule
\end{tabular}
\end{center}
\noindent{\small $^{\ast}$ Value of $k$ inferred by us from $H(49)=240$ and $H(48)=236$, not read from the source.}

The bound $246$ stood from 2014 until the last days of August 2026. It has been
formalised in Lean~4, assuming only the Bombieri--Vinogradov theorem
\cite{Axiom26}. The bounds from $240$ to $186$ appeared within one week.
They rest on Stadlmann's idea of enlarging the support of the sieve weights
using equidistribution estimates for moduli with a large smooth factor. The
Lean development accompanying the bound $186$ is conditional on three stated
inputs (two exponential-sum bounds and a family of numerical integral bounds)
\cite{OAI26}. An independent re-run of that repository, dated 7 and 8 November
2026, reproduced the build and the numerical certificate and describes the
result as ``conditionally supported, with the computation independently
reproduced; not yet settled''; the two exponential-sum bounds remain axioms
\cite{NM26}. A further document, dated 18 September 2026 and naming an AI
system as its author, claims $\DHL[39,2]$ and hence $182$; by its own abstract
its Lean development verifies the last implication ``under explicit numerical
and finite-field premises'' \cite{Sha26}. We had missed it in earlier
versions. (These details were read by an earlier run of this project; on
12 November the file was too large for our reader and direct download was
refused, so we have not re-checked them since.) As of 7 September 2026 a standard reference still listed $246$ as
the best fully established value \cite{MW}, and the page read the same on 12
November. On the date of this version we found no report of complete
independent verification of any of the 2026 claims and no claimed bound below
$182$. We therefore keep two columns throughout: \emph{established} ($246$)
and \emph{claimed} ($182$).

\subsection*{Why \texorpdfstring{$2$}{2} is a different problem}
Under the generalised Elliott--Halberstam conjecture the sieve gives $H_1\le6$,
and Polymath8b showed that $6$ cannot be improved by purely sieve-theoretic
arguments: this is Selberg's parity obstruction \cite{Pol14b}. Lowering the
bound (Track~A below) and proving the conjecture (Track~B) are therefore
different tasks. No amount of success on Track~A reaches $2$.

\section{What this project has verified for itself}

\subsection{The combinatorial half, checked by the Lean kernel}

Our Lean development is written in core Lean~4, without Mathlib. It defines
primality, the twin prime conjecture, the predicate \texttt{BoundedGaps B}
(infinitely many pairs of distinct primes at distance at most $B$),
admissibility, and $\DHL[k,2]$. It contains no \texttt{sorry}, and the only
axioms it uses are \texttt{propext} and \texttt{Quot.sound}. The statement
$\DHL[k,2]$ is always a hypothesis, never an axiom.

\begin{theorem}[\lbl{Lean-checked}]\label{thm:lean}
\begin{enumerate}
\item The twin prime conjecture is equivalent to \texttt{BoundedGaps 2}.
\item If a finite Boolean check succeeds on a list $\mathcal H$ (it misses a
class modulo every $m$ with $2\le m\le|\mathcal H|$), then $\mathcal H$ is
admissible. The moduli $p>|\mathcal H|$ are handled by a formal pigeonhole
argument.
\item For each pair $(k,D)$ in Table~\ref{tab:tuples}, for
$(k,D)=(56,278)$, $(57,282)$, $(58,288)$, $(59,300)$, $(60,304)$, $(61,310)$,
$(62,320)$, $(63,324)$, $(64,330)$, and for $(102,578)$ and $(105,600)$, there is an explicit admissible $k$-tuple in $[0,D]$, and
$\DHL[k,2]$ implies \texttt{BoundedGaps D}. In particular $\DHL[2,2]$ implies
the twin prime conjecture.
\item $\DHL[k,2]$ implies $\DHL[k',2]$ for every $k'\ge k$
(\texttt{DHL.mono}: delete shifts from the tuple). So a proof of $\DHL[j,2]$
with $j\le k$ gives \texttt{BoundedGaps D} for each listed pair $(k,D)$.
\end{enumerate}
\end{theorem}

\begin{table}[h]
\caption{Pairs $(k,D)$ with a kernel-checked admissible $k$-tuple of diameter
$D$. The tuples were found by our own search; each $D$ agrees with the value
of $H(k)$ in published tables, whose optimality we have not re-proved.}
\label{tab:tuples}
\begin{tabular}{@{}rr@{\qquad}rr@{\qquad}rr@{\qquad}rr@{\qquad}rr@{\qquad}rr@{}}
\toprule
$k$&$D$&$k$&$D$&$k$&$D$&$k$&$D$&$k$&$D$&$k$&$D$\\
\midrule
2&2&11&36&20&80&29&130&38&176&47&226\\
3&6&12&42&21&84&30&136&39&182&48&236\\
4&8&13&48&22&90&31&140&40&186&49&240\\
5&12&14&50&23&94&32&146&41&188&50&246\\
6&16&15&56&24&100&33&152&42&196&51&252\\
7&20&16&60&25&110&34&156&43&200&52&254\\
8&26&17&66&26&114&35&158&44&210&53&264\\
9&30&18&70&27&120&36&162&45&212&54&270\\
10&32&19&76&28&126&37&168&46&216&55&272\\
\bottomrule
\end{tabular}
\end{table}

Theorem~\ref{thm:lean} is elementary. Its value is that implication
\eqref{eq:reduction} is now a checked interface: any future proof of
$\DHL[k,2]$, by anyone, yields the corresponding gap bound with nothing left
to trust on the combinatorial side.

\subsection{The sieve constant, with exact certificates}

Let $\R_k=\{t\in[0,1]^k:\ t_1+\dots+t_k\le1\}$ and, for square-integrable $F$
supported on $\R_k$,
\[
I_k(F)=\int_{\R_k}F^2,\qquad
J_k^{(m)}(F)=\int\left(\int F\,dt_m\right)^{2}\prod_{i\ne m}dt_i,\qquad
M_k=\sup_F\frac{\sum_{m=1}^{k}J_k^{(m)}(F)}{I_k(F)}.
\]
Maynard's theorem \cite[Prop.~4.2]{May15} states that if the primes have level
of distribution $\theta$ and $\theta M_k/2>1$ then $\DHL[k,2]$ holds. With
Bombieri--Vinogradov ($\theta<1/2$) the condition is $M_k>4$.

We take $F=\sum_{a+2b\le d}c_{a,b}(1-P_1)^aP_2^b$ with $P_j=\sum_it_i^j$, as
in \cite[\S8]{May15}. All integrals reduce to
\[
\int_{\R_k}(1-P_1)^AP_2^B\,dt=\frac{A!}{(k+A+2B)!}\,G_k(B),\qquad
\sum_{B\ge0}G_k(B)\frac{x^B}{B!}=\left(\sum_{b\ge0}\frac{(2b)!}{b!}x^b\right)^{k},
\]
which follows from the Dirichlet integral. The two quadratic forms are then
exact rational matrices. We compute an approximate top eigenvector in
floating point, round it to rationals, and evaluate the Rayleigh quotient in
exact arithmetic; the last step is a rigorous lower bound whatever happened
before it.

\begin{theorem}[\lbl{Computed}, exact rational arithmetic]\label{thm:Mk0}
With $F$ a polynomial in $1-P_1$ and $P_2$ only,
\[
M_{105}\ge4.00206976,\qquad M_{102}\ge4.00271343,\qquad M_{5}\ge2.00410541 .
\]
\end{theorem}

The first value reproduces Maynard's $M_{105}>4$ with $d=11$ ($42$
coefficients). The third gives $\DHL[5,2]$ and $H_1\le12$ under the
Elliott--Halberstam conjecture, again reproducing \cite{May15}. For
$d=11,15,19$ the least $k$ at which this family certifies $M_k>4$ is
$105,102,102$, and at $k=101$ it gives only $3.99931$ ($d=19$): polynomials in
$P_1,P_2$ alone stop at $k=102$.

\subsection{All symmetric test functions}
To go further we allow every symmetric polynomial. Since $P_1,\dots,P_k$
generate the symmetric polynomials, the functions
\[
(1-P_1)^a\,P^{b}:=(1-P_1)^a\prod_{j\ge2}P_j^{b_j},\qquad a+w(b)\le d,\quad w(b)=\sum_j jb_j,
\]
span all symmetric polynomials of degree at most $d$ (for $k\ge d$). The
integrals have a closed form that we have not seen stated in this way
(\lbl{Proved here}; possibly folklore).

\begin{lemma}\label{lem:egf}
Let $X_1,\dots,X_k$ be independent exponential random variables of mean $1$
and $E_k(b)=\mathbb E\prod_jP_j(X)^{b_j}$. Then
\[
\int_{\R_k}(1-P_1)^AP^{b}\,dt=\frac{A!}{(k+A+w(b))!}\,E_k(b),
\qquad
\sum_b E_k(b)\frac{s^b}{b!}=\left(\sum_{c}w(c)!\,\frac{s^c}{c!}\right)^{k},
\]
where $s^b=\prod_js_j^{b_j}$ and $b!=\prod_jb_j!$. In particular each
$E_k(b)$ is a positive integer, computed exactly by multiplying
exponential generating functions.
\end{lemma}

\begin{proof}
For a monomial, the Dirichlet integral gives
$\int_{\R_k}(1-P_1)^At^\alpha\,dt=A!\,\alpha!/(k+A+|\alpha|)!$, and
$\alpha!=\mathbb EX^\alpha$. Summing over the monomials of $P^b$, all of
degree $w(b)$, gives the first identity. For the second,
$\mathbb E\exp(\sum_js_jP_j(X))=\left(\mathbb E\exp(\sum_js_jX_1^j)\right)^k$
by independence, and
$\mathbb E\exp(\sum_js_jX_1^j)=\sum_c\mathbb E X_1^{w(c)}\,s^c/c!=\sum_cw(c)!\,s^c/c!$,
as formal power series.
\end{proof}

The $J$-functional reduces to the same lemma in $k-1$ variables: with
$P_j=P_j'+t_k^j$,
\[
\int_0^{1-P_1'}(1-P_1)^aP^b\,dt_k=\sum_{e\le b}\binom{b}{e}
\frac{a!\,w(e)!}{(a+w(e)+1)!}\,(1-P_1')^{a+w(e)+1}P'^{\,b-e}.
\]
So both Gram matrices are exact rational matrices. We find an approximate top
eigenvector by inverse iteration with a fixed shift above
$\frac{k}{k-1}\log k$, which can only converge to the top eigenvalue. We take
its floating-point midpoints as exact dyadic rationals and evaluate the
Rayleigh quotient exactly. The working precision of the eigenvector matters
only for speed. A separate script, \texttt{code/verify\_cert.py}, rebuilds
both Gram matrices from scratch and re-evaluates each stored coefficient
vector in exact arithmetic.

\begin{theorem}[\lbl{Computed}, exact rational arithmetic]\label{thm:Mk}
\[
M_{105}\ge4.17427,\quad M_{64}\ge4.00363,\quad M_{58}\ge4.00927,\quad
M_{56}\ge4.00519,\quad M_{55}\ge4.00516 .
\]
The first two use all symmetric polynomials of degree $11$ and $12$
($195$ and $272$ coefficients). The others use the subspace spanned by
$(1-P_1)^aP^b$ with $b$ supported on even $j$, as in \cite[\S7]{Pol14b}, of
degree $16$, $18$ and $20$ ($307$, $471$ and $707$ coefficients).
\end{theorem}

\begin{center}
\begin{tabular}{@{}lccc@{}}
\toprule
Basis & degree & coefficients & least $k$ with certified $M_k>4$\\
\midrule
$1-P_1$, $P_2$ & 19 & 110 & 102\\
all symmetric & 12 & 272 & 64\\
even power sums & 16 & 307 & 58\\
even power sums & 18 & 471 & 56\\
even power sums & 20 & 707 & 55$^{\dagger}$\\
\bottomrule
\end{tabular}
\end{center}
\noindent{\small $^{\dagger}$ Only this value of $k$ was run at this degree; smaller $k$ were not tried.}

\begin{corollary}[end-to-end within this project]\label{cor:578}
Assuming the Bombieri--Vinogradov theorem and \cite[Prop.~4.2]{May15}, both
\lbl{Literature}, Theorems~\ref{thm:Mk} and~\ref{thm:lean} give
$H_1\le272$.
\end{corollary}

Polymath8b reached $M_{54}>4.00238$ with even-part polynomials of degree $23$
\cite[Thm.~3.9]{Pol14b}, and hence $H_1\le270$. Our pipeline is one step of
$k$ short of that; the degree-$23$ computation for $k=54$ is running.
Corollary~\ref{cor:578} reproduces known results and is weaker than the
established $246$. It is the position of our own pipeline on the ladder
$600\to270\to246\to240\to212\to186\to182$, which we intend to climb rung by
rung.

\subsection{What the plain method cannot do}
Polymath8b proved $M_k\le\frac{k}{k-1}\log k$ \cite{Pol14b}. Hence
(\lbl{Computed} from that bound) $M_k>4$ forces $k\ge51$, and on the plain
simplex $\DHL[k,2]$ requires a level of distribution
$\theta>2(k-1)/(k\log k)$:
\begin{center}
\begin{tabular}{@{}rcccccccc@{}}
\toprule
$k$ & 50 & 45 & 40 & 35 & 30 & 20 & 10 & 5\\
\midrule
$\theta>$ & 0.501 & 0.514 & 0.529 & 0.547 & 0.568 & 0.634 & 0.782 & 0.994\\
\bottomrule
\end{tabular}
\end{center}
The 2026 results go below $k=50$ by changing the support of $F$, not by
proving a uniform level of distribution of this size.

\section{Results of our own: rigidity of the pigeonhole step}

Every bound above ends with the same step: if the weighted average number of
primes among $n+h_1,\dots,n+h_k$ exceeds $1$, two of them are prime, and they
are at most $h_k-h_1$ apart. A natural hope is to do better with unequal
weights. Fix $D$, and give each position a weight $c_i\ge0$. Call a set of
positions \emph{$D$-separated} if any two of its shifts differ by more than
$D$, and let $\alpha_D(c)$ be the largest total weight of a $D$-separated set.
If no two primes among the $n+h_i$ are within $D$, then
$\sum_ic_i\mathbf 1_{n+h_i\text{ prime}}\le\alpha_D(c)$. Writing
\[
M_k[c]=\sup_F\frac{\sum_{m}c_mJ_k^{(m)}(F)}{I_k(F)},
\]
the sieve therefore yields $H_1\le D$ as soon as
$\tfrac{\theta}{2}M_k[c]>\alpha_D(c)$. For $c=(1,\dots,1)$ and $D=h_k-h_1$
this is the usual criterion. Could a cleverer $c$ certify some $D<h_k-h_1$
that no sub-tuple certifies?

\begin{proposition}[\lbl{Proved here}; possibly folklore]\label{prop:rigid}
Let $\omega$ be the largest number of shifts $h_i$ contained in an interval of
length $D$. Then for every $c\ge0$, not identically zero,
\[
\frac{M_k[c]}{\alpha_D(c)}\le M_\omega .
\]
So the weighted criterion never certifies more than the standard criterion
applied to the best sub-tuple of diameter at most $D$.
\end{proposition}

\begin{proof}
Let $G$ be the graph on $\{1,\dots,k\}$ joining $i\ne j$ when
$|h_i-h_j|\le D$. Its stable sets are the $D$-separated sets and its cliques
are the sets of shifts lying in an interval of length $D$. Normalise
$\alpha_D(c)=1$, so $c$ lies in
\[P=\{c\ge0:\ c(S)\le1\text{ for every stable set }S\}.\] The graph $G$ is a
unit interval graph, hence perfect, and so is its complement by the perfect
graph theorem \cite{Lov72}. For a perfect graph the polytope cut out by the
clique inequalities is the convex hull of the indicator vectors of its stable
sets \cite{Chv75}. Applying this to the complement of $G$ gives
\[P=\operatorname{conv}\{\mathbf 1_K:\ K\text{ a clique of }G\}.\]

The map $c\mapsto M_k[c]$ is a supremum of linear functionals, hence convex,
and therefore $M_k[c]\le\max_KM_k[\mathbf 1_K]$. It remains to show
$M_k[\mathbf 1_K]\le M_{j}$ for $|K|=j$; this suffices because $M_j$ is
non-decreasing in $j$ (adjoin a variable and concentrate $F$ near its zero
set). Write $t=(s,u)$ with $s$ the coordinates in $K$ and $u$ the others, put
$\sigma(u)=1-\sum u\ge0$, and let $F_u=F(\cdot,u)$, a function on the scaled
simplex $\sigma(u)\R_j$. For $m\in K$ we have
$J_k^{(m)}(F)=\int J_j^{(m)}(F_u)\,du$, and $I_k(F)=\int I_j(F_u)\,du$.
Rescaling $F_u$ to $\R_j$ multiplies $I_j$ by $\sigma^{j}$ and each
$J_j^{(m)}$ by $\sigma^{j+1}$, so
\[
\sum_{m\in K}J_j^{(m)}(F_u)\ \le\ \sigma(u)\,M_j\,I_j(F_u)\ \le\ M_j\,I_j(F_u).
\]
Integrating over $u$ gives the claim.
\end{proof}

\begin{remark}[\lbl{Computed}, consistency check for $k=3$]
With unrestricted polynomials of degree $6$ on $\R_3$ we find the lower bounds
$M_3[(1,1,1)]\ge1.6464$, $M_3[(1,1,0)]\ge1.3579$ and
$M_3[(\tfrac12,1,\tfrac12)]\ge1.1604$, against
$M_2=1/(1-W(1/e))=1.3859\ldots$\,. The half-weighted functional is far below
$M_2$, as the proposition requires.
\end{remark}

\subsection*{The same rigidity on the enlarged simplex}
The $\varepsilon$-enlarged simplex of Polymath8b \cite[Thm.~3.12]{Pol14b}
takes $F$
supported on $(1+\varepsilon)\R_k$ and restricts the outer integral of each
$J^{(m)}$ to $(1-\varepsilon)\R_{k-1}$. Call the resulting constant
$M_{k,\varepsilon}$. The criterion $M_{k,\varepsilon}>2/\theta$ is valid in
two cases (\lbl{Literature}): under $\mathrm{EH}[\theta]$ when
$1+\varepsilon<1/\theta$, and under $\mathrm{GEH}[\theta]$ when
$\varepsilon<1/(k-1)$. With Bombieri--Vinogradov the first case allows every
$\varepsilon<1$, and Polymath8b's $H_1\le246$ uses $M_{50,1/25}>4.0043$
\cite[Thm.~3.13]{Pol14b}. Unequal weights are no more useful there.

\begin{proposition}[\lbl{Proved here}; possibly folklore]\label{prop:rigid-eps}
For $0\le\varepsilon<1$ and every weight $c\ge0$, $c\ne0$,
$M_{k,\varepsilon}[c]/\alpha_D(c)\le M_{\omega,\varepsilon}$, with equality
for the indicator of a largest clique. Since the conditions on $\varepsilon$
above do not get harder as $k$ decreases, the weighted criterion never
certifies more than the standard one for the best sub-tuple of diameter at
most $D$.
\end{proposition}

\begin{proof}
As for Proposition~\ref{prop:rigid}, it suffices to show
$\sum_{m\in K}J_{k,\varepsilon}^{(m)}(F)\le M_{j,\varepsilon}I_k(F)$ for a
clique $K$ of size $j$, together with $M_{j,\varepsilon}\le M_{\omega,\varepsilon}$.
Fix the coordinates $u$ outside $K$, with $|u|\le1-\varepsilon$ (other slices
do not meet the $J$-region). The slice $F_u$ is supported on
$(1+\varepsilon-|u|)\R_j$, and its $J$-region is
$(1-\varepsilon-|u|)\R_{j-1}$. Put $\mu=(1+\varepsilon-|u|)/(1+\varepsilon)\le1$
and $H(x)=F_u(\mu x)$, which is supported on $(1+\varepsilon)\R_j$. Then
$I_j(F_u)=\mu^jI_j(H)$, and each slice $J$-term equals $\mu^{j+1}$ times the
integral of $(\int H\,dx_m)^2$ over $|x'|\le(1-\varepsilon-|u|)/\mu$. Since
\[
\frac{(1-\varepsilon-|u|)(1+\varepsilon)}{1+\varepsilon-|u|}\le1-\varepsilon
\iff -|u|(1+\varepsilon)\le-|u|(1-\varepsilon),
\]
that region lies inside $(1-\varepsilon)\R_{j-1}$, so the slice contributes at
most $\mu^{j+1}M_{j,\varepsilon}I_j(H)=\mu M_{j,\varepsilon}I_j(F_u)$. For
$j=1$ use Cauchy--Schwarz, which gives $M_{1,\varepsilon}=1+\varepsilon$.
Integrate over $u$ and use $\mu\le1$. Monotonicity in $j$ follows by
embedding: take $F(s,u)=H(s/\mu_\delta)\mathbf 1_{[0,\delta]^{k-j}}(u)$ with
$\mu_\delta=1-(k-j)\delta/(1+\varepsilon)$ and let $\delta\to0$.
\end{proof}

Slicing naively gives each slice its own parameter
$\varepsilon'=\varepsilon/(1-|u|)\ge\varepsilon$, which would only bound the
weighted constant by $\sup_{\varepsilon'\ge\varepsilon}M_{j,\varepsilon'}$, a
weaker statement. The single rescaling by $\mu$ above avoids that loss.

\begin{remark}[\lbl{Computed}, $k=3$]
With unrestricted polynomials (degree $5$ for $k=3$, degree $10$ for $k=2$,
cut off at $|t|\le\rho$ for $\rho\in\{1-\varepsilon,1,1+\varepsilon\}$):
\begin{center}
\begin{tabular}{@{}ccccc@{}}
\toprule
$\varepsilon$ & $M_{2,\varepsilon}\ge$ & $M_{3,\varepsilon}[(1,1,0)]\ge$ & $M_{3,\varepsilon}[(\frac12,1,\frac12)]\ge$ & $M_{3,\varepsilon}\ge$\\
\midrule
$0$ & 1.385933 & 1.349844 & 1.160121 & 1.646427\\
$1/10$ & 1.518567 & 1.482143 & 1.264698 & 1.793492\\
$1/5$ & 1.627161 & 1.585674 & 1.341926 & 1.882282\\
$1/3$ & 1.717104 & 1.663743 & 1.413309 & 1.917001\\
\bottomrule
\end{tabular}
\end{center}
Both weighted columns stay below $M_{2,\varepsilon}$, as the proposition
requires.
\end{remark}

\begin{remark}[scope]
Neither proposition covers more general polytope supports
\cite[\S3]{Pol14b}, Stadlmann's hybrid supports \cite{Sta26}, $J$-regions
that depend on $m$, or arguments that are not linear in the indicators of
primality. Those are \lbl{Open}.
\end{remark}

\section{Results of our own: how far the \texorpdfstring{$\varepsilon$}{epsilon}-enlargement alone can go}

For the plain simplex the bound $M_k\le\frac{k}{k-1}\log k$ shows where the
method stops ($k\ge51$). This section starts the same analysis for the
$\varepsilon$-enlarged simplex. For $0<\tau\le L$ put
\[
M_k(L,\tau)=\sup_F\frac{\sum_{m=1}^k\int_{\tau\R_{k-1}}\left(\int_0^\infty F\,dt_m\right)^2\prod_{i\ne m}dt_i}{\int F^2},
\]
the supremum over square-integrable $F\ne0$ supported on $L\R_k$. Then
$M_{k,\varepsilon}=M_k(1+\varepsilon,1-\varepsilon)$ and $M_k=M_k(1,1)$. The
quantity $M_k(L,\tau)$ is non-decreasing in $L$ and in $\tau$, satisfies
$M_k(cL,c\tau)=c\,M_k(L,\tau)$, and is non-decreasing in $k$ by the embedding
used in the proof of Proposition~\ref{prop:rigid-eps}. Write
$S=t_1+\dots+t_k$ and call the index $m$ \emph{active} at $t$ if
$S-t_m\le\tau$.

\begin{lemma}[\lbl{Proved here}; it is {\cite[Lemma~6.1]{Pol14b}} with the region of the marginals kept]\label{lem:dual}
Let $\psi_1,\dots,\psi_k$ be measurable functions on $L\R_k$, positive almost
everywhere, such that for almost every $t'\in\tau\R_{k-1}$ and every $m$
\[
\int_0^{L-|t'|}\frac{du}{\psi_m(t',u)}\le1,
\]
where $(t',u)$ is the point whose $m$-th coordinate is $u$ and whose other
coordinates are $t'$. Then
\[
M_k(L,\tau)\le\operatorname*{ess\,sup}_{t\in L\R_k}\ \sum_{m\text{ active at }t}\psi_m(t).
\]
\end{lemma}

\begin{proof}
For $t'\in\tau\R_{k-1}$ the Cauchy--Schwarz inequality gives
\[
\left(\int F\,dt_m\right)^2\le\int\frac{dt_m}{\psi_m}\cdot\int\psi_mF^2\,dt_m\le\int\psi_mF^2\,dt_m .
\]
Integrating over $t'\in\tau\R_{k-1}$ gives the integral of $\psi_mF^2$ over
the set where $m$ is active. Sum over $m$.
\end{proof}

\begin{proposition}[\lbl{Proved here}; elementary, possibly folklore]\label{prop:eps-first}
For $k\ge2$ and $0\le\varepsilon<1$,
\[
M_{k,\varepsilon}\le(1+\varepsilon)M_k\le(1+\varepsilon)\frac{k}{k-1}\log k .
\]
Conversely, for $0\le\varepsilon\le M_k/(4k)$,
\[
M_{k,\varepsilon}\ge(1+\varepsilon)M_k-\frac{64k^2}{M_k}\,\varepsilon^2 .
\]
In particular $\lim_{\varepsilon\to0^+}(M_{k,\varepsilon}-M_k)/\varepsilon=M_k$.
\end{proposition}

\begin{proof}
By monotonicity in $\tau$ and scaling,
$M_k(1+\varepsilon,1-\varepsilon)\le M_k(1+\varepsilon,1+\varepsilon)=(1+\varepsilon)M_k$,
and $M_k\le\frac{k}{k-1}\log k$ is \cite[Cor.~6.4]{Pol14b}.

For the lower bound, scaling gives $M_{k,\varepsilon}=(1+\varepsilon)M_k(1,\rho)$
with $\rho=(1-\varepsilon)/(1+\varepsilon)$. Put $\delta=k(1-\rho)$ and let $F$
on $\R_k$ satisfy $\sum_mJ_k^{(m)}(F)\ge(M_k-\eta)I_k(F)$. Split $F=F_1+F_2$
with $F_2=F\mathbf 1_{S>\rho}$, and let $x=I_k(F_2)/I_k(F)$. Two facts follow
from the Cauchy--Schwarz inequality on segments of length at most $1-\rho$.
First, a fiber of $F_2$ has length at most $1-\rho$, so
$\sum_mJ_k^{(m)}(F_2)\le\delta\,I_k(F_2)$. Second, a fiber over $t'$ with
$|t'|>\rho$ lies in $\{S>\rho\}$ and has length at most $1-\rho$, so
restricting the marginals of $F$ to $\rho\R_{k-1}$ loses at most
$\delta\,I_k(F_2)$. By the triangle inequality for
$F\mapsto(\int F\,dt_m)_{m\le k}$ and the first fact,
\[
M_k-\eta\le\left(\sqrt{M_k(1-x)}+\sqrt{\delta x}\right)^2\le M_k(1-x)+2\sqrt{M_k\delta x}+\delta x .
\]
This is a quadratic inequality in $\sqrt x$. If $\delta\le M_k/2$ it gives
$x\le(16\delta+8\eta)/M_k$. By the second fact the quotient of $F$ for
$M_k(1,\rho)$ is at least $M_k-\eta-\delta x$. Letting $\eta\to0$ we get
$M_k(1,\rho)\ge M_k-16\delta^2/M_k$. Finally
$\delta=2k\varepsilon/(1+\varepsilon)$, and $\delta\le M_k/2$ when
$\varepsilon\le M_k/(4k)$.
\end{proof}

So for fixed $k$ the factor $1+\varepsilon$ is exact to first order, and no
dual argument can improve it there. Any saving is an effect of large $k$: it
needs $\varepsilon$ of size about $1/k$ or more, where the shell
$\rho<\sigma\le1$ carries most of the mass. The first inequality already
restricts the parameter: with Bombieri--Vinogradov the criterion
$M_{k,\varepsilon}>4$ forces
$\varepsilon>\varepsilon_0(k)=\frac{4(k-1)}{k\log k}-1$ (\lbl{Computed}):
\begin{center}
\begin{tabular}{@{}rccccccc@{}}
\toprule
$k$ & 49 & 48 & 45 & 40 & 39 & 35 & 30\\
\midrule
$\varepsilon_0(k)$ & 0.0068 & 0.0117 & 0.0274 & 0.0572 & 0.0638 & 0.0929 & 0.1369\\
\bottomrule
\end{tabular}
\end{center}

For larger $\varepsilon$ we use the lemma with a weight that changes shape
where the indices start to become inactive.

\begin{theorem}[\lbl{Proved here} for the inequality; \lbl{Computed}, interval arithmetic, for the numerical statement]\label{thm:floor}
Let $k\ge2$, $0<\tau<s_1\le L$ and $a\ge s_1$. Then
\[
M_k(L,\tau)\le\sup_{0\le\sigma\le\tau}W(\sigma),\qquad
W(\sigma)=\frac{ka}{k-1}\log\left(1+\frac{(k-1)(s_1-\sigma)}{a-\sigma}\right)+L-s_1+\sigma\log\frac{L-\sigma}{s_1-\sigma}.
\]
Consequently $M_{k,\varepsilon}<3.99$ for every $k\le29$ and every
$0\le\varepsilon<1$. In particular the criterion of \cite[Thm.~3.12]{Pol14b}
on the $\varepsilon$-enlarged simplex cannot give $\DHL[k,2]$ for any
$k\le29$ from the Bombieri--Vinogradov theorem, nor from any level of
distribution $\theta\le0.501$, whatever the value of $\varepsilon$.
\end{theorem}

\begin{proof}
Fix $\Lambda>0$, put $\lambda=\Lambda/(ka)$ and take the two-zone weight
\[
\psi_m(t)=\lambda\,(a-S+kt_m)\quad(S\le s_1),\qquad
\psi_m(t)=\Lambda\,\frac{t_m}{S}\quad(S>s_1).
\]
It is positive wherever $t_m>0$. Where $S\le s_1\le a$ every term is
non-negative, so the sum over the active indices is at most the sum over all
indices, which is $\lambda(k(a-S)+kS)=\Lambda$. Where $S>s_1$ the sum over the
active indices is at most $\Lambda\sum_mt_m/S=\Lambda$. On the fiber over $t'$
with $|t'|=\sigma\le\tau$ we have $S=\sigma+u$, and
\[
\int_0^{s_1-\sigma}\frac{du}{\lambda(a-\sigma+(k-1)u)}+\int_{s_1-\sigma}^{L-\sigma}\frac{(\sigma+u)\,du}{\Lambda u}=\frac{W(\sigma)}{\Lambda}.
\]
Lemma~\ref{lem:dual} with $\Lambda=\sup W$ gives the inequality.

For the numerical statement it suffices to treat $k=29$. For
$\varepsilon\le143/1000$, Proposition~\ref{prop:eps-first} gives
$M_{29,\varepsilon}\le1.143\cdot\frac{29}{28}\log29<3.9863$. The interval
$[143/1000,1]$ is cut into $857$ intervals $[e,e']$ of length $1/1000$. For
each one the file \texttt{code/cert\_floor\_k29.json} lists rational $a$ and
$s_1$ with $\sup W<3.99$ for $(L,\tau)=(1+e',1-e)$; by monotonicity in $L$ and
$\tau$ this covers every $\varepsilon\in[e,e']$. The supremum over $\sigma$ is
bounded rigorously: the first logarithm is non-increasing in $\sigma$ (because
$a\ge s_1$) and the last term is non-decreasing, so on a piece $[x,y]$ we bound
$W$ by the first term at $x$ plus the last at $y$, in ball arithmetic at
rational points, bisecting until the bound is below $3.99$. The script is
\texttt{code/dual\_eps\_cert.py}; a second script,
\texttt{code/check\_floor\_cert.py}, re-checks the certificate with a
different interval library. The largest value of $\sup W$ over the $857$
intervals is $3.9772$ (floating point).
\end{proof}

\begin{remark}[what was known, and a caution]
Polymath8b state the uniform bound $M_{k,\varepsilon}\le\frac{k}{k-1}\log(2k-1)$
for $0\le\varepsilon<1$ \cite[Prop.~6.5]{Pol14b}, and remark that the bounds
for $M_{k,\varepsilon}$ ``do not appear to be as sharp''. That bound is below
$4$ exactly for $k\le23$, so Theorem~\ref{thm:floor} moves this floor from $24$
to $30$. It does not use their bound. We note, with caution, that we could not
complete the last step of the proof of \cite[Prop.~6.5]{Pol14b} as we read it:
the weights $1-S+kt_m$ are summed, and for $S>1$ their sum over the active
indices can exceed $k$ (at $t=(s,0,\dots,0)$ with $1<s\le1+\varepsilon$ only
the first index is active and its weight is $1+(k-1)s$). Our reconstruction
gives the constant $\frac{k}{k-1}+\varepsilon$ in place of $\frac{k}{k-1}$. We
read that proof only through a text extraction of the arXiv version and may be
misreading it; this is \lbl{Open} until checked against the published text. A
brief search (the paper \cite{Pol14b} and the Polymath wiki page on the
variational problem) found no sharper published upper bound for
$M_{k,\varepsilon}$; the search was not exhaustive, and
Theorem~\ref{thm:floor} has not been reviewed by anyone outside this project.
\end{remark}

\begin{remark}[\lbl{Computed}, floating point only: how far this weight goes]
The best value of $\sup W$ over $(a,s_1)$, maximised over $\varepsilon$, is
about $3.975$ for $k=29$, $4.004$ for $k=30$, $4.257$ for $k=40$ and $4.459$
for $k=50$, in each case near $\varepsilon=0.24$ to $0.28$; the bound of
\cite[Prop.~6.5]{Pol14b} is $4.187$, $4.218$, $4.482$ and $4.689$. For small
$\varepsilon$ the two-zone weight does not beat the scaling bound: at $k=50$
and $\varepsilon=1/25$ both give $4.15$ ($4.1458$ against $4.1515$), while
Polymath8b's certified lower bound is $4.0043$. The true threshold is presumably
close to $k=50$, so the theorem is far from sharp. Script:
\texttt{code/dual\_eps.py remark15}.
\end{remark}

\begin{remark}[a weight that did not work]
We also tried the weight $\psi_m=\frac{\beta}{n}(a-S_A+nt_m)$, where $n$ is
the number of active indices at $t$ and $S_A$ the sum of the active
coordinates. Its sum over the active indices is exactly $\beta a$ everywhere,
so the bound is $a$ times the largest fiber integral. Evaluating the fiber
integral exactly on the fibers over $t'=(x,\dots,x,0,\dots,0)$ (\lbl{Computed},
floating point, $k=50$, $\varepsilon=1/25$, every $a\in[1.05,3]$ on a grid of
step $0.01$, script \texttt{code/dual\_eps.py remark16}) shows that this
weight cannot give less than $4.178$, which is worse than the $4.1515$ from
scaling. The binding fiber has a single non-zero coordinate $x\approx2\varepsilon$:
there the zero coordinates become inactive while $n$ is still large, and the
normalisation by $n$ wastes the slack. A dual bound that is sharper for small
$\varepsilon$ must treat the small coordinates near such sparse points
differently. This is \lbl{Open}. (An earlier draft of this remark, version~0.3
as first circulated, reported a weaker obstruction, $4.146$, from fibers with
about ten non-zero coordinates; the exact fiber evaluation above supersedes it.)
\end{remark}

\begin{remark}[scope]
Theorem~\ref{thm:floor} concerns the $\varepsilon$-enlarged simplex only. It
says nothing about the 2026 claims, which use other supports and new
equidistribution estimates. A floor for those is the Problem in the next
section.
\end{remark}

\section{Programme}

\subsection*{Track A: lowering the bound}
\begin{enumerate}
\item \emph{Climb the ladder.} The certified pipeline stands at $k=55$
($H_1\le272$). Next: $M_{54}>4$ ($H_1\le270$; even power sums of degree
$23$), then add the
$\varepsilon$-enlarged simplex ($k=50$, $246$), then Stadlmann's region
($k=49$, $240$).
\item \emph{Floor of the method (our main original direction).} For the plain
simplex the dual bound $M_k\le\frac{k}{k-1}\log k$ shows where the method
stops ($k\ge51$). For the $\varepsilon$-enlarged simplex
Theorem~\ref{thm:floor} gives $k\ge30$, which is still far from the expected
truth near $50$; Section~4 explains what a sharper weight must do. No upper
bound of this kind is known to us for the 2026 hybrid supports. A rigorous
dual certificate there would say whether $182$ is close to the limit of the
present inputs or far from it, and which input (which Type~I/II/III range) is
binding.
\item \emph{Which rung is cheap.} From Table~\ref{tab:tuples}, the gain per
unit of $k$ below the claimed frontier is $182\to176\to168\to162\to158$
for $k=39,\dots,35$: the steps $39\to38$ and $38\to37$ are worth the most.
\end{enumerate}

\begin{problem}[\lbl{Open}]
Give a rigorous upper bound for the variational quantity of \cite{Sta26} as a
function of $k$ and of the admissible exponents, and deduce the least $k$ the
hybrid method can reach with the currently proved estimates.
\end{problem}

\subsection*{Track B: the parity barrier}
No sieve of the above kind can distinguish integers with an even number of
prime factors from those with an odd number. Known ways around the barrier
all inject bilinear or structural information. We will study, one per
session, what each would need in order to say something about twin primes:
Chen's theorem \cite{Che73}; Bombieri's asymptotic sieve \cite{Bom76}; the
parity-breaking sieve of Friedlander and Iwaniec \cite{FI98}; Heath-Brown's
theorem that Siegel zeros would imply the conjecture \cite{HB83}; the
function-field proof of Sawin and Shusterman \cite{SS22}; and the
entropy-based results on correlations of multiplicative functions
\cite{Tao16}.

\begin{problem}[\lbl{Open}, to be made precise in Lean]
State, as a single explicit hypothesis on bilinear sums over $n(n+2)$, the
weakest known input that implies the twin prime conjecture, and formalise the
implication.
\end{problem}

\subsection*{Rules of this paper}
A numerical bound enters the ``established by this project'' column only with
an exact certificate and, for the combinatorial part, a Lean proof. Ideas are
logged with their outcome, including failures. Nothing here is peer reviewed.

\begin{thebibliography}{99}
\bibitem{Axiom26} E.~Chen, S.~Hariharan, K.~Lau, B.~Mehta, A.~A.~Swaminathan,
J.~Thorner, E.~Xie, \emph{Lean formalization of bounded gaps between primes}
(blueprint, PrimeGapsLib), 2026. \url{https://primegaps.axiommath.ai/}
\bibitem{Bom76} E.~Bombieri, \emph{The asymptotic sieve}, 1976.
\bibitem{Cha26} F.~Charton, L.~Hong, K.~Lau, K.~Ono, G.~Remy, H.~C.~Siu,
A.~A.~Swaminathan, J.~Thorner, Y.~Xie, \emph{A new bound for small gaps
between primes}, preliminary draft, 3 September 2026.
\url{https://primegaps.axiommath.ai/bgp212.pdf}
\bibitem{Che73} J.~R.~Chen, \emph{On the representation of a larger even
integer as the sum of a prime and the product of at most two primes}, 1973.
\bibitem{Chv75} V.~Chv\'atal, \emph{On certain polytopes associated with
graphs}, J. Combin. Theory Ser. B, 1975.
\bibitem{FI98} J.~Friedlander, H.~Iwaniec, \emph{The polynomial $X^2+Y^4$
captures its primes}, Ann. of Math., 1998.
\bibitem{HB83} D.~R.~Heath-Brown, \emph{Prime twins and Siegel zeros}, Proc.
London Math. Soc., 1983.
\bibitem{Kin26} S.~Kintali, \emph{Improved bounded gaps between primes},
preprint, 1 September 2026.
\bibitem{Lov72} L.~Lov\'asz, \emph{Normal hypergraphs and the perfect graph
conjecture}, Discrete Math., 1972.
\bibitem{May15} J.~Maynard, \emph{Small gaps between primes}, Ann. of Math.
181 (2015), 383--413.
\bibitem{MW} E.~W.~Weisstein, \emph{Bounded gaps between primes}, MathWorld.
\url{https://mathworld.wolfram.com/BoundedGapsBetweenPrimes.html}
\bibitem{NM26} Nervous Machine, \emph{DEBRIEF-math-primegaps}, a re-run of
the repository \texttt{openai/PrimeGaps186} by a third party (``Scorecard, not
a verification''), status dated 7 November 2026.
\url{https://github.com/Nervous-Machine/DEBRIEF-math-primegaps}
\bibitem{OAI26} OpenAI, \emph{Improved short gaps between primes}, preprint
dated 30 August 2026, with the repository \texttt{openai/PrimeGaps186}.
\bibitem{Pol14b} D.~H.~J.~Polymath, \emph{Variants of the Selberg sieve, and
bounded intervals containing many primes}, Res. Math. Sci. 1 (2014).
\bibitem{Sha26} GPT-6 Astra (OpenAI), prompted by R.~Sharifi, \emph{A
fixed-profile proof for prime gaps of at most 182}, dated 18 September 2026
(first version 4 September 2026).
\url{https://www.math.ucla.edu/~sharifi/AI/prime_gaps_182_expository_bundle.pdf}
\bibitem{SS22} W.~Sawin, M.~Shusterman, \emph{On the Chowla and twin primes
conjectures over $\mathbb F_q[T]$}, Ann. of Math., 2022.
\bibitem{SY26} Z.~Song, S.~Yue, \emph{Bounded gaps between primes: an upper
bound of 236}, IACR ePrint 2026/1893.
\bibitem{Sta26} J.~Stadlmann, \emph{Bounded gaps between primes},
arXiv:2608.31126, 31 August 2026.
\bibitem{Tao16} T.~Tao, \emph{The logarithmically averaged Chowla and Elliott
conjectures for two-point correlations}, Forum Math. Pi, 2016.
\bibitem{Zha14} Y.~Zhang, \emph{Bounded gaps between primes}, Ann. of Math.
179 (2014), 1121--1174.
\end{thebibliography}

\end{document}
